\documentclass[../main.tex]{subfiles}
\begin{document}

Дадлагын хүрээнд хэрэглэгчийн байршуулсан баримтуудад тулгуурлан асуултад
хариулдаг, эх сурвалжаа ишэлдэг, монгол хэлийг бүрэн дэмждэг RAG чатбот
системийг төлөвлөлтөөс эхлээд ажиллаж буй бүтээгдэхүүн хүртэл бүтээж
дуусгав. Систем нь нэг Fastify API, PostgreSQL + pgvector сан, вэб болон
мобайл гэсэн хоёр клиентээс бүрдэх бөгөөд бүртгэл, нэвтрэлт, олон
хэрэглэгчийн тусгаарлалттайгаар бодит төхөөрөмжүүд дээр туршигдсан.

\section{Төлөвлөгөөний биелэлт}

Удирдагчаар батлуулсан төлөвлөгөөний (Хүснэгт~\ref{tab:internship-plan})
12 ажил бүгд хэрэгжсэн. Гүйцэтгэлийн явцад гарсан хазайлтуудыг нуулгүй
дурдвал:

\begin{itemize}
\item Арын боловсруулалтад анх төлөвлөсөн BullMQ дарааллын системийн оронд
Node.js-ийн \texttt{setImmediate} ашигласан --- нэг серверийн MVP-д тусдаа
Redis хамаарал нэмэх нь илүүц гэж үндэслэсэн (4.1 хэсэг);
\item хариултыг үсэг үсгээр цацах SSE (streaming) төлөвлөгөөнд байсан ч
MVP хүрээнээс хасаж, төлөв шалгах polling-оор шийдсэн;
\item хайлтын \texttt{threshold}-ын анхдагч утга онолын 0.7-оос 0.1 болж
буусан нь хэл дамнасан хайлтын түр шийдэл бөгөөд үндсэн шийдэл нь
асуултын орчуулга гэж баримтжуулагдсан (6.2 хэсэг).
\end{itemize}

Нөгөө талд төлөвлөгөөнөөс давсан ажил ч хийгдсэн: мобайл аппликейшн анхны
хүрээнд байгаагүй ч бүрэн үйлдэлтэй клиент болтлоо хөгжиж, вэбтэй ижил
боломжуудтай болсон. Дадлага нь их сургуулиар үзсэн өгөгдлийн сан,
алгоритмын онолын мэдлэгийг бодит системийн шаардлагатай холбож, судалгаа
$\rightarrow$ зохиомж $\rightarrow$ хэрэгжүүлэлт $\rightarrow$ туршилт
гэсэн инженерийн бүрэн мөчлөгийг нэг төсөл дээр туулах боломж олгов.

\section{Байгууллагад өгөх санал}

Бүтээсэн системийг байгууллагын үйл ажиллагаанд дараах байдлаар нэвтрүүлэх
боломжтой гэж үзэж байна. Дотоод журам, гарын авлага, тайлан зэрэг байнга
лавлагддаг баримтуудыг системд индексжүүлснээр ажилтнууд «энэ асуудлыг
аль журмаар зохицуулдаг вэ?» маягийн асуултын хариуг эх сурвалжийн ишлэлтэй
нь хормын дотор авах дотоод мэдлэгийн туслах болгон ашиглаж болно. Мөн
үйлчилгээний нийтлэг асуултуудын баримтад тулгуурлан хэрэглэгчийн
асуулт-хариултын (help center) эхний шатлалыг автоматжуулах чиглэл бий.
Нэвтрүүлэхийн өмнө гурван зүйлийг шийдвэрлэхийг зөвлөж байна: системийг
байнгын сервер дээр байршуулж нөөцлөлтийн (backup) горим тогтоох; ашиглалтын
хэмжээнээс хамааран Gemini API-ийн төлбөрт түвшинд шилжих эсэхийг үнэлэх
(үнэгүй түвшний квотын хүлээлт 6.5 хэсгийн хэмжилтээр гол сааталт);
баримтын хандалтын эрхийн бодлогыг тодорхойлох (одоогийн систем
хэрэглэгч бүрийн баримтыг бүрэн тусгаарладаг тул хэлтсийн түвшний
хуваалцах эрх нэмэх шаардлага гарч болзошгүй).

\section{Цаашдын хөгжүүлэлт}

Системийг цааш хөгжүүлэх тодорхой чиглэлүүд:

\begin{itemize}
\item \textbf{Асуултын орчуулга (query translation).} Асуултыг баримтын
хэл рүү орчуулж хайдаг болгосноор threshold-ыг онолын түвшинд нь эргүүлэн
өсгөж, хэл дамнасан хайлтын чанарыг үндсээр нь сайжруулах;
\item \textbf{Production байршуулалт ба CI/CD.} Системийг үүлэн сервер дээр
HTTPS-тэйгээр байршуулж, шинэчлэлтийг автоматаар хүргэх CI/CD дамжлага
болон өдөр тутмын өгөгдлийн нөөцлөлт тохируулах;
\item \textbf{Мобайл түгээлт.} EAS build-ээр суулгах боломжтой APK/IPA
гаргаж, Expo Go-оос хараат бус болгох;
\item \textbf{Скан баримтын дэмжлэг (OCR).} Зураг хэлбэрээр хадгалагдсан
PDF-ээс текст ялгах давхарга нэмж, дэмжигдэх баримтын хүрээг тэлэх;
\item \textbf{PDF доторх хүснэгт, зургийн нарийвчлалтай боловсруулалт.}
Одоогийн систем PDF-ээс зөвхөн урсгал текст ялгадаг тул хүснэгтийн мөр,
баганын бүтэц алдагдаж, зураг болон диаграм огт орхигддог. Хүснэгтийг
бүтэцтэй хэлбэрээр нь (мөр, баганын харгалзааг хадгалан) ялгаж
индексжүүлэх, зураг диаграмыг multimodal загвараар тайлбар текст болгон
хөрвүүлж хамт хадгалах чадвар нэмбэл тоон үзүүлэлт, график агуулгатай
холбоотой асуултын хариултын нарийвчлал мэдэгдэхүйц сайжирна;
\item \textbf{Хариултын streaming.} Хариултыг бүрэн үүсч дуустал хүлээхийн
оронд үсэг үсгээр цацаж, хүлээлтийн мэдрэмжийг багасгах.
\end{itemize}

Ерөнхийд нь дүгнэвэл, дадлага нь тавьсан зорилтоо бүрэн биелүүлж, онолын
судалгаанаас бодит хэрэглээ хүртэлх замыг бүтэн туулсан, цаашид байгууллагын
орчинд нэвтрүүлэх боломжтой систем болон дадлагажигчийн хувьд production
түвшний хөгжүүлэлтийн иж бүрэн туршлага гэсэн хоёр талын үр дүнтэй болов.

\end{document}
